\noindent \textit{\textcolor{red!60!black}{
Major Comment 3.\\
Another substantial concern is the lack of clarity on what exactly "leverage risk" is, is it new, what it would predict, and if those differentially predictions are really demonstrated empirically. The authors mention debt overhang and also models of uncertainty, but those seem like they might generate completely opposing predictions. If there are costs of financial distress then it seems likely that firms would want to build-up cash, rather than engage in equity payouts, if they are concerned about future uncertainty (aka an increased probability of needing to make larger debt payments to avoid distress). 
}}

\textit{\textcolor{red!60!black}{
On the other hand, if it is about debt overhang then perhaps as the authors indicate firms could want to increase equity payouts. If debt overhang is the mechanism though, then there does already exist some empirical evidence in that space (ex. Bennett et al. 2018 and Wittry 2018 as recent examples) and it isn't clear in the way things are written now that leverage risk is empirically distinguishable from debt overhang in this setting. As the authors note in this setting conditional on the abrogation being overturned there could be a 69\% increase in bond payments for exposed firms just after perhaps the most severe economic downturn in U.S. history. It seems entirely plausible this has a substantial effect on the probability of bankruptcy, and therefore debt overhang. Firms always have some "leverage risk" since the future value of the asset side of the balance sheet is always uncertain and that could generate debt overhang if firms are sufficiently levered. What then are the differential predictions when there is "leverage risk" coming from uncertainty in the future liability and how do the authors show evidence of that? If there aren't differential predictions, then it might be more appropriate to frame this paper as additional evidence on the effects of debt overhang on corporate investment and then compare/contrast much more with the existing literature/evidence in that space.
}}

\bigskip

\noindent Thank you for this comment. We agree, based on feedback from the other reviewers as well, that the term ``leverage risk'' was confusing. We have removed it throughout the paper, including from the title, which is now ``Debt Overhang During Policy Uncertainty: The Case of Gold Clauses in the 1930s.''

As you correctly note, firms always face some degree of leverage risk due to uncertainty in the future value of assets, which can generate debt overhang if firms are sufficiently levered. Our setting differs in that the source of uncertainty arises on the liability side of the balance sheet---specifically, from the potential revaluation of gold clause obligations. We view the predictions from gold clause exposure as a direct application of debt overhang. To clarify this point, we added Section 5.1, which frames our analysis within the dynamic debt overhang model of \cite{Hennessy2004}. As we write in the manuscript:

\begin{quote}
``For firms with long-term debt, \cite{Hennessy2004} shows that capital structure creates a wedge $r_{j,t}$ between a firm's \textit{average} $Q$ and its \textit{marginal} $q$... The gold clause litigation provides a natural setting to apply this framework. A ruling against the abrogation would have raised the claim of debtholders on the firm's assets and increased the probability of default. Accordingly, the legal uncertainty about the abrogation must have increased the expected recovery value $R_{j,t}$ and thus the debt overhang wedge $r_{j,t}$.''
\end{quote}

We interpret 1933 and 1934 as years when this wedge as a function of a firm's gold clause exposure. 

We have also expanded our discussion of the broader debt overhang literature in Section 2, explicitly positioning our contribution relative to existing empirical work. As we now write in the introduction: ``Our results complement the empirical literature on debt overhang that exploits plausibly exogenous variation in leverage to study firm decisions (e.g., \cite{Giroud2012}, \cite{Wittry2018}, \cite{Bennett2025}). With respect to this literature, our empirical setting has a unique feature: the abrogation created uncertainty about future leverage increases that ultimately did not materialize.'' This unique feature---that the liability increase never materialized---connects our work to the literature on firm-level uncertainty, which we now discuss more explicitly following R2's suggestions.

Regarding your concern about opposing predictions from debt overhang versus precautionary motives: we find that firms with higher gold exposure raised equity payouts rather than accumulating cash and reducing leverage. This pattern is consistent with \cite{Hennessy2004}'s prediction that shareholders of distressed firms optimally extract value through dividends, and harder to reconcile with precautionary behavior that would predict lower payouts and cash hoarding. Together, this evidence suggests that debt overhang is the more likely operative mechanism.
